\documentclass[12pt, a4paper]{article}

% Pakiety podstawowe i obsługa języka polskiego
\usepackage[utf8]{inputenc}
\usepackage[T1]{fontenc}
\usepackage[polish]{babel}

% Pakiety do formatowania
\usepackage{lmodern}
\usepackage{graphicx}
\usepackage{amsmath}
\usepackage{amssymb}
\usepackage{ulem} % do przekreślania tekstu (\sout)
\usepackage{booktabs} % do ładnych tabel (\toprule, \midrule, \bottomrule)
\usepackage{hyperref} % do linków
\usepackage{xcolor} % do kolorowania składni
\usepackage{parskip}

% Pakiet do kodów źródłowych
\usepackage{listings}
\lstset{
    basicstyle=\ttfamily\small,
    backgroundcolor=\color{gray!10},
    frame=single,
    breaklines=true,
    captionpos=b,
    literate={ą}{{\k{a}}}1 {ć}{{\'c}}1 {ę}{{\k{e}}}1 {ł}{{\l{}}}1 {ń}{{\'n}}1 {ó}{{\'o}}1 {ś}{{\'s}}1 {ź}{{\'z}}1 {ż}{{\.z}}1 {Ą}{{\k{A}}}1 {Ć}{{\'C}}1 {Ę}{{\k{E}}}1 {Ł}{{\L{}}}1 {Ń}{{\'N}}1 {Ó}{{\'O}}1 {Ś}{{\'S}}1 {Ź}{{\'Z}}1 {Ż}{{\.Z}}1
}

% Konfiguracja linków
\hypersetup{
    colorlinks=true,
    linkcolor=blue,
    filecolor=magenta,      
    urlcolor=blue,
}

\title{Markdown, Colab i GitHub}
\author{}
\date{}

\begin{document}

\maketitle

Na początek poznamy trzy narzędzia, które rozwiązują trzy różne problemy:

\begin{table}[h]
\centering
\begin{tabular}{lp{10cm}}
\toprule
\textbf{Narzędzie} & \textbf{Do czego służy?} \\
\midrule
Markdown & do prostego zapisywania uporządkowanego tekstu, dokumentacji i raportów \\
Google Colab & do łączenia tekstu, kodu, obliczeń i wyników w jednym notebooku \\
GitHub & do przechowywania projektu, śledzenia zmian i współpracy \\
\bottomrule
\end{tabular}
\end{table}

Te trzy elementy bardzo dobrze ze sobą współpracują. Podstawą wszędzie jest przygotowanie opisu w kodzie Markdown. Do tego dochodzić będą obliczenia w Colabie, a następnie przechowanie źródła i wyniki w repozytorium GitHub.

\section{Markdown}

\subsection{Co to jest Markdown?}
Markdown to prosty język znaczników służący do formatowania zwykłego tekstu.
Najczęściej zapisujemy go w plikach z rozszerzeniem: \texttt{.md}

Przykłady:
\begin{itemize}
    \item \texttt{README.md}
    \item \texttt{raport.md}
    \item \texttt{notatki.md}
\end{itemize}

Plik Markdown jest zwykłym plikiem tekstowym. Można go otworzyć praktycznie w każdym edytorze. Specjalne znaki, takie jak \texttt{\#}, \texttt{*}, \texttt{-} czy \texttt{\$}, informują program wyświetlający dokument, jak ma wyglądać tekst.

\subsection{Dlaczego Markdown jest tak istotny?}
Za każdym razem rozmawiając z chatem AI tak naprawdę rozmawiamy z plikiem Markdown. AI potrafi interpretować składnię Markdown i generować raporty w tym formacie. 

Dlatego też jak skopiujesz output to ma on całą masę znaczników (np. \texttt{\#}, \texttt{*}, \texttt{-}, czy \texttt{\$}), które są potrzebne do poprawnego renderowania.

Github (archiwum programistyczne) zawiera opisy również w Markdown. Podobnie Google Colab (środowisko programistyczne pythona) potrafią renderować Markdown. Dzięki temu możemy przygotować raport, który będzie wyglądał dobrze zarówno w przeglądarce, jak i w notebooku.

Warto od razu rozróżnić dwie rzeczy:
\begin{itemize}
    \item \textbf{źródło Markdown} --- tekst, który wpisujemy do pliku,
    \item \textbf{wyrenderowany dokument} --- wygląd tego tekstu po interpretacji składni Markdown.
\end{itemize}

Na GitHubie możemy przełączać się między widokiem źródła i wyrenderowanym dokumentem. 

Podobnie edytory kodu jak VS Code, PyCharm czy Jupyter Notebook potrafią renderować Markdown w osobnym oknie podglądu.

\subsection{Zanim zaczniemy pisać w Markdown zobacz przykładowy plik Markdown.}
Raport na temat ciała na sprężynie: 
\begin{itemize}
    \item \glqq Ciało na sprężynie\grqq{}. 
\end{itemize}


\section{Nagłówki}
Nagłówki budują strukturę dokumentu. Im więcej znaków \texttt{\#}, tym niższy poziom nagłówka.

\textbf{Źródło:}
\begin{lstlisting}[language=bash]
# Tytuł dokumentu
## Głowna sekcja
### Podsekcja
\end{lstlisting}

\textbf{Efekt:} \\
{\Large \textbf{Tytuł dokumentu}} \\
{\large \textbf{Główna sekcja}} \\
{\normalsize \textbf{Podsekcja}}

W praktyce jeden dokument powinien mieć jeden główny nagłówek \texttt{\#}. Dalsze części dokumentu organizujemy za pomocą \texttt{\#\#}, \texttt{\#\#\#} itd.

\section{Pogrubienie, kursywa i przekreślenie}

\textbf{Źródło:}
\begin{lstlisting}
**tekst pogrubiony**
*tekst pochylony*
~~tekst przekreslony~~
***tekst pogrubiony i pochylony***
\end{lstlisting}

\textbf{Efekt:}\\
\textbf{tekst pogrubiony}\\
\textit{tekst pochylony}\\
\sout{tekst przekreślony}\\
\textbf{\textit{tekst pogrubiony i pochylony}}

Wyróżnień warto używać oszczędnie. Jeśli wszystko jest pogrubione, nic nie jest naprawdę wyróżnione.

\section{Akapit i nowa linia}
Nowy akapit tworzymy przez pozostawienie pustej linii.

\textbf{Źródło:}
\begin{lstlisting}
To jest pierwszy akapit.

To jest drugi akapit.
\end{lstlisting}

\textbf{Efekt:}\\
To jest pierwszy akapit.

To jest drugi akapit.

To ważna zasada: w Markdown pojedyncze naciśnięcie \texttt{Enter} nie zawsze oznacza nowy akapit.

\section{Listy punktowane}
\textbf{Źródło:}
\begin{lstlisting}
- Markdown
- Colab
- GitHub
\end{lstlisting}

\textbf{Efekt:}
\begin{itemize}
    \item Markdown
    \item Colab
    \item GitHub
\end{itemize}

Listy można zagnieżdżać:
\begin{lstlisting}
- GitHub
  - repozytorium
  - Issues
  - historia zmian
- Colab
  - tekst
  - kod
  - wyniki
\end{lstlisting}
\textbf{Efekt:}
\begin{itemize}
    \item GitHub
    \begin{itemize}
        \item repozytorium
        \item Issues
        \item historia zmian
    \end{itemize}
    \item Colab
    \begin{itemize}
        \item tekst
        \item kod
        \item wyniki
    \end{itemize}
\end{itemize}
\newpage
\section{Listy numerowane}
\textbf{Źródło:}
\begin{lstlisting}
1. przygotuj dane,
2. wykonaj obliczenia,
3. zapisz wynik,
4. opisz wynik.
\end{lstlisting}

\textbf{Efekt:}
\begin{enumerate}
    \item przygotuj dane,
    \item wykonaj obliczenia,
    \item zapisz wynik,
    \item opisz wynik.
\end{enumerate}
Listy numerowane są szczególnie przydatne przy opisie procedury lub instrukcji.

\section{Checklisty}
GitHub obsługuje listy zadań.

\textbf{Źródło:}
\begin{lstlisting}
- [x] utworzono notebook
- [x] wykonano obliczenia
- [ ] sprawdzono wynik
\end{lstlisting}

\textbf{Efekt:}
\begin{itemize}
    \item[$\boxtimes$] utworzono notebook
    \item[$\boxtimes$] wykonano obliczenia
    \item[$\square$] sprawdzono wynik
\end{itemize}
Checklisty bardzo dobrze nadają się do krótkich list kontrolnych w dokumentacji i GitHub Issues.
\newpage
\section{Linki}
\textbf{Źródło:}
\begin{lstlisting}
[GitHub](https://github.com/)
\end{lstlisting}

\textbf{Efekt:}\\
\href{https://github.com/}{GitHub}

Zamiast wklejać długi adres do strony, warto umieścić go pod czytelną nazwą.

\section{Obrazy}
Obraz osadzamy podobnie jak link, ale na początku dodajemy znak \texttt{!}.

\textbf{Źródło:}
\begin{lstlisting}
![Wykres zaleznosci y od x](wykres.png)
\end{lstlisting}

Jeżeli plik \texttt{wykres.png} znajduje się w tym samym folderze co dokument Markdown, GitHub wyświetli obraz w miejscu tej instrukcji. Jeżeli plik znajduje się w podfolderze \texttt{files}, używamy ścieżki:
\begin{lstlisting}
![Wykres zaleznosci y od x](files/wykres.png)
\end{lstlisting}

\textbf{Efekt:}
\begin{figure}[h]
    \centering
    \includegraphics[width=1.0\textwidth]{files/wykres.png}
\end{figure}

Tekst w nawiasach kwadratowych jest opisem obrazu. Warto go uzupełniać --- pomaga zrozumieć zawartość także wtedy, gdy obraz nie może zostać wyświetlony. 

W projektach przechowywanych w repozytorium zwykle wygodniej używać ścieżek względnych niż pełnych adresów internetowych.

\section{Tabele}
\textbf{Źródło:}
\begin{lstlisting}
| Narzedzie | Typ pliku | Zastosowanie |
| --- | --- | --- |
| Markdown | .md | dokumentacja |
| Colab | .ipynb | obliczenia |
| GitHub | repozytorium | historia projektu |
\end{lstlisting}

\textbf{Efekt:}
\begin{table}[h]
\centering
\begin{tabular}{lll}
\toprule
\textbf{Narzędzie} & \textbf{Typ pliku} & \textbf{Zastosowanie} \\
\midrule
Markdown & \texttt{.md} & dokumentacja \\
Colab & \texttt{.ipynb} & obliczenia \\
GitHub & repozytorium & historia projektu \\
\bottomrule
\end{tabular}
\end{table}

Tabele są dobre do krótkich zestawień. Jeśli w komórkach pojawiają się długie akapity tekstu, zwykle lepiej wrócić do zwykłych sekcji i list.

\section{Kod w tekście}
Nazwy plików, komendy i bardzo krótkie fragmenty kodu warto oznaczać pojedynczym backtickiem.

\textbf{Źródło:}\\
\texttt{Uruchom plik `analiza.py` i sprawdź wartość zmiennej `wynik`.}

\textbf{Efekt:}\\
Uruchom plik \texttt{analiza.py} i sprawdź wartość zmiennej \texttt{wynik}.

\section{Bloki kodu}
Większy fragment kodu umieszczamy w osobnym bloku.

\textbf{Źródło:}
\begin{lstlisting}
```python
x = 5
y = x**2
print(y)
```
\end{lstlisting}

Podanie nazwy języka, np. \texttt{python}, \texttt{bash}, \texttt{html} czy \texttt{json}, pozwala GitHubowi kolorować składnię.

W raporcie technicznym dobrze jest pokazać krótki fragment kodu, który jest istotny dla opisywanego wyniku. Nie trzeba kopiować do raportu całego programu.

\section{Cytaty}

\textbf{Źródło:}
\begin{verbatim}
> Wynik należy zawsze sprawdzić przed jego wykorzystaniem.
\end{verbatim}

\textbf{Efekt:}
\begin{quote}
Wynik należy zawsze sprawdzić przed jego wykorzystaniem.
\end{quote}

Cytaty można wykorzystać także do krótkiego wyróżnienia komentarza lub uwagi.

\section{Matematyka}
GitHub potrafi renderować wzory matematyczne. Krótki wzór umieszczamy wewnątrz pojedynczych znaków dolara.

\textbf{Źródło:}
\begin{lstlisting}
Energia spoczynkowa dana jest wzorem $E=mc^2$.    
\end{lstlisting}
\textbf{Efekt:}

Energia spoczynkowa dana jest wzorem $E=mc^2$. Dłuższy wzór umieszczamy w osobnym bloku.

\textbf{Źródło:}
\begin{lstlisting}
$$
\sum_{k=1}^{n} k^2 = \frac{n(n+1)(2n+1)}{6}
$$
\end{lstlisting}

\textbf{Efekt:}
$$\sum_{k=1}^{n} k^2 = \frac{n(n+1)(2n+1)}{6}$$
Macierze zapisujemy wielowierszowo:

\textbf{Źródło:}
\begin{lstlisting}
$$
A= 
\begin{pmatrix} 
1 & 2 \\ 
3 & 4 \\ 
\end{pmatrix}
$$
\end{lstlisting}

\textbf{Efekt:}
$$A= \begin{pmatrix} 1 & 2 \\ 3 & 4 \\ \end{pmatrix}$$

\section{Zrzut ekranu i transkrypcja do Markdown}
Czasem informacja istnieje tylko jako obraz: zrzut ekranu programu, zdjęcie tabeli, komunikat błędu albo fragment zeskanowanego dokumentu. AI może pomóc przepisać zawartość obrazu do tekstu. Taki proces nazywamy tutaj \textbf{transkrypcją}. 

Przykładowy workflow:
\begin{enumerate}
\item zachowujemy oryginalny obraz,
\item prosimy AI o \textbf{przepisanie jego zawartości do kodu markdown}
\item zapisujemy wynik w pliku tekstowym lub Markdown,
\item porównujemy transkrypcję z obrazem,
\item poprawiamy ewentualne błędy.
\end{enumerate}Największej ostrożności wymagają:
\begin{itemize}
\item liczby,
\item wzory,
\item nazwy plików,
\item adresy,
\item fragmenty kodu,
\item komunikaty błędów.
\end{itemize}
Transkrypcja wygenerowana przez AI jest \textbf{wersją roboczą}, dopóki człowiek jej nie sprawdzi.

\section{Poprawne formatowanie matematyki w Markdown}
W materiałach tego kursu polecamy stosowanie prostego zestawu reguł formatowania matematyki przygotowanego specjalnie tak, aby pliki dobrze renderowały się na GitHubie: 
\begin{itemize}
    \item \texttt{hints.md}. 
\end{itemize}
 Pracując z chatem AI i domagając się kodu markdowna, warto załączać ten plik jako kontekst, aby AI od razu generowało poprawny kod w Markdown. 
 
 Różne systemy: przeglądarkowe chaty AI, Edytor Visual Studio, Colab i GitHub, mogą różnie renderować wzory matematyczne!

Częstym błędem jest używanie w kodzie Markdown znaków \verb|\[]| i \verb|\]| zamiast \verb|$ $| do oznaczania wzorów matematycznych. Wtedy GitHub nie renderuje wzoru, a Colab renderuje go poprawnie.

\section{Dobre materiały o Markdown}
\begin{itemize}
\item \href{https://docs.github.com/en/get-started/writing-on-github/getting-started-with-writing-and-formatting-on-github/basic-writing-and-formatting-syntax}{GitHub Docs --- Basic writing and formatting syntax}
\item \href{https://docs.github.com/en/get-started/writing-on-github/working-with-advanced-formatting/writing-mathematical-expressions}{GitHub Docs --- Writing mathematical expressions}
\item \href{https://www.markdownguide.org/basic-syntax/}{Markdown Guide --- Basic Syntax}
\item \href{https://www.markdownguide.org/cheat-sheet/}{Markdown Guide --- Cheat Sheet}
\end{itemize}

\end{document}
